\chapter{Further actual-angle reductions from exact source entry}
\label{component:pe56fb394b522}
\noindent\textbf{Component }\path{pe56fb394b522}\par
{\small\noindent Source: \path{R13/baseline/baseline/baseline/baseline/advances/source/angle_reductions.tex}\par}
{\small\noindent SHA256: \path{e56fb394b52263918d985b9437e107edea077b0ea105beebe65fca8b1c306250}\par}
{\small\noindent Inventory use: Actual directed weak II/IU/UI/UU exclusions and both-strong UU five-event entry\par}
\medskip
\section[Further actual-angle reductions from exact source entry]{Further actual-angle reductions from exact source entry}
\label{pe56fb394b522:sec:c2-angle-reductions}

All angles in this section are actual source contacts. They are strictly
positive finite field elements; ``infinite'' below means unbounded in
the coefficient field, not a literal projective endpoint. Write
$z=c^2=u_c^2$, $y=a^2z=v_c^2$, and
$A_c=1+zr$, $S=A_c+a^2zt$. The exact source identities used are
\[
 n\ge0\ \Longleftrightarrow\ (b+r)A_cS-r\Delta\ge0,
 \qquad
 N=(1+z)(b+r)h+a^2\{g+j+rt(1+y)\}>0.
\]

\begin{theorem}[The ordinary--ordinary source orbit is already closed]
If both actual C2 half-angles are positive bounded units, then the
source belongs either to the two-weak-complementary-event all-tests
exit or to the five-consecutive-strong-event all-tests exit. No upper
bound on an original event coefficient is required.
\end{theorem}
\begin{proof}
Here $a,z$ are units. The positive gaps give $0<h<1/z$ and
$0<t<1+zr$, so $h,t,b=h+a^2t$ are bounded. From
$g=b-1+r+zh>0$ we obtain $1<b+r+zh\le(1+z)b+r$;
thus $b+r$ is a positive unit. The positive expansion of $\Delta$
makes it a bounded unit too. The original $A$ formula has unit
numerator and bounded denominator $h$, so $A$ is strong.

The exact numerator for $n$, after substituting $t=(b-h)/a^2$,
satisfies the polynomial specialization
\[
 \left[(b+r)A_cS-r\Delta\right]_{r=1}
     =-z(a^2+h)(b+1)(z+1).
\]
If $1-r$ were infinitesimal, taking residues would contradict
$n\ge0$: the right side is strictly negative. Hence $1-r$ is a
unit. The original $m$ numerator is now a unit and $g\Delta$ is
bounded, so $m$ is strong. Apply the exact C2 reflection; it preserves
ordinary actual angles and exchanges $m,D$. Thus $D$ is strong.

If either $C$ or $\alpha$ is weak, its original rational formula
implies $N$ infinitesimal: all the multiplying denominator factors
are bounded and $z,a,b+r,\Delta$ have bounded inverses. The positive
identity for $N$ then forces $h,j$ infinitesimal. Consequently
$t=A_c-j$ and $1-zh$ are units. Both original complementary formulas
therefore give $C,\alpha$ infinitesimal. The two-weak theorem applies.
Otherwise $C,\alpha$ are both strong; with $m,A,D$ strong, the
five-consecutive-event theorem applies. These cases exhaust the source.
\end{proof}

\paragraph{Some directed one-weak angle types are empty.}
Orient an exactly-one-weak source so that $C$ is weak and $\alpha$
strong. In addition to the ordinary--ordinary exclusion above, its
actual angle types cannot be $(\mathrm{unit},\mathrm{Inf})$,
$(\mathrm{Inf},\mathrm{unit})$ or $(\mathrm{Inf},\mathrm{Inf})$.
Here $\mathrm{Inf}$ denotes positive infinitesimal.

For $(\mathrm{unit},\mathrm{Inf})$, $z$ is a unit and $a^2$ is
infinitesimal. Again $t,h,b,\Delta$ are bounded, $\Delta$ is a unit,
and $b+r$ is a unit. Weak $C$ implies $N/a^2$ infinitesimal.
The positive identity forces $j$ and $rt$ infinitesimal. Since
$t=A_c-j$ is a unit, $r$ is infinitesimal. Then $g>0$ makes
$b+zh$ a unit, and $b=h+a^2t$ makes $h$ a unit. The first positive
summand of $N$ is consequently a unit, contradicting $N/a^2$
infinitesimal.

For $(\mathrm{Inf},\mathrm{unit})$, bounded $b$ is a unit by the
inherited source-entry inequality. Explicitly,
\[
 xC(b-h)=\frac{A_c(xN)}{2\Delta},\qquad
 xN=(1+x)(b+r)(xh+k^2)-x(1+k^2)(1-r)(b-h).
\]
The last numerator has residue $k_0^2(b_0+r_0)>0$, independently
of any lower bound on $g$. Thus $xC(b-h)$ is a unit and $C$ is strong. For unbounded
$b$, (SR1) gives $C\asymp1/L$ with bounded $L$, again strong.
This exhausts $b$ and excludes weak $C$ in this angle orientation.

For $(\mathrm{Inf},\mathrm{Inf})$, both $z,y$ are infinitesimal
and $t<A_c$ is bounded. Put $H=zh$, $L=z(b-h)=yt$ and $B=zb=H+L$.
Here $0<H<1$, $L$ is infinitesimal, and $B$ is bounded. If $H$ is
a unit then $b$ is unbounded; direct division of the original
$\Delta$ by $b$ gives a positive unit of residue $1+B_0$.
The positive first summand of $N$ consequently gives
$C\ge c_0H/L$ for some positive real $c_0$, contradicting weak $C$.
Thus weak $C$ forces $H,B$ infinitesimal. The exact ratio is then
\[
 \frac\alpha C=
 \frac{z^2(b+r)(b-h)}{A_c(1-H)}
 =\frac{(B+zr)L}{A_c(1-H)},
\]
which is infinitesimal. Therefore $\alpha$ is weak as well, contrary
to the selected one-weak orientation.

After these exclusions, an exactly-one-weak source oriented with weak
$C$ has only the following candidate actual types:
\begin{gather*}
 (\mathrm{Inf},\mathrm{unbounded}),\quad
 (\mathrm{unit},\mathrm{unbounded}),\quad
 (\mathrm{unbounded},\mathrm{Inf}),\\
 (\mathrm{unbounded},\mathrm{unit}),\quad
 (\mathrm{unbounded},\mathrm{unbounded}).
\end{gather*}
The second type is the reflection of the single-angle domain and
is now fully covered by the homogeneous unbounded-$b$ theorem and the
other single-angle admissions. The four other listed directed types
remain conditional source leaves; no nonemptiness is inferred from
their listing.

\paragraph{Other strong-event angle types are not all empty.}
For $x>0$ infinitesimal, the exact positive chart
\[
 c=\sqrt x,\quad a=1,\quad b=1/x,\quad r=1/4,\quad t=1/2
\]
has $h=1/x-1/2$, $g=1/x+1/4-x/2$, $j=1/2+x/4$ and all
remaining chart gaps positive. Its original coefficients are
\begin{align*}
 n&=\frac{8-5x^2}{16(1+x)^2},&
 m&=\frac{3x^2(3x+4)}{16(1+x)(4+x-2x^2)},\\
 C=\alpha&=\frac{4-x}{8x},&
 A&=\frac{x(x+4)(3x+4)}{32(2-x)(1+x)^2},\\
 D&=\frac{(2-x)(3x+4)}{16x(1+x)}.
\end{align*}
Every coefficient is positive. The unchanged exact contact is
\[
 (u_c,v_c;W,X,Y,Z)
 =(\sqrt x,\sqrt x;1/x,-\sqrt x/4,1/\sqrt x-\sqrt x/2,1).
\]
Original source stationarity, positive transverse state inertia and
Laurent convexity prove physicality, including projective limits.
Here $C,\alpha,D$ are unbounded, $m\sim3x^2/16$, $A\sim x/4$,
and $n\to1/2$. Thus a genuine strong-complementary source can have
two infinitesimal actual angles and weak $m,A$. This is an obstruction
to deleting that angle orbit, not a counterexample to tensor SOS.
